% Práctica 05 | Restauración del respaldo.
\section{Restauración del respaldo}
\subsection{Entorno de ejecución}
La restauración de \texttt{ejemplo.backup} se realizó el 26 de octubre de
2026 a las 22:15, hora de la Ciudad de México. Se utilizó PostgreSQL 16.15
en un contenedor de la imagen oficial \texttt{postgres:16}, denominado
\texttt{nipgg-p05-restauracion}, sin puertos publicados y con la red
deshabilitada.

Se creó la base \texttt{ejemplo} a partir de \texttt{template0}. Esta base
se destinó exclusivamente a la restauración; las pruebas de restricciones
se ejecutaron por separado. Antes de cargar el respaldo, se comprobó que
el esquema \texttt{public} no contenía tablas.

\subsection{Identificación del respaldo}
\begingroup\small\renewcommand{\arraystretch}{1.4}
\begin{tabularx}{\linewidth}{@{}p{.30\linewidth}J@{}}
\toprule\textbf{Propiedad} & \textbf{Valor observado}\\\midrule
Archivo & \texttt{ejemplo.backup}\\
Formato & \texttt{CUSTOM}, versión 1.15; compresión gzip\\
Versión de origen & PostgreSQL 16.3; \texttt{pg\_dump} 16.3\\
Fecha del respaldo & 24 de marzo de 2025, 00:57:23 UTC\\
Índice de contenido & 310 entradas\\
Extensión incluida & \texttt{adminpack}\\\bottomrule
\end{tabularx}\endgroup

Se eligió PostgreSQL 16 para mantener la compatibilidad con la extensión
\texttt{adminpack} incluida en el archivo. La identificación del insumo y
la comprobación de la base vacía se muestran en la figura~\ref{fig:inspeccion}.

\subsection{Procedimiento de restauración}
Se copió el respaldo al contenedor, se examinó su índice con
\texttt{pg\_restore --list} y se ejecutó el siguiente comando:
\begin{lstlisting}[style=sqlacademico]
docker exec nipgg-p05-restauracion pg_restore \
  --verbose --clean --if-exists --no-acl --no-owner \
  --exit-on-error --single-transaction \
  -U postgres -d ejemplo /tmp/ejemplo.backup
\end{lstlisting}

Las opciones \texttt{--clean --if-exists} eliminan los objetos anteriores
solo si existen; \texttt{--no-acl --no-owner} omite los privilegios y
propietarios del entorno de origen. La transacción única permite revertir
la operación completa si ocurre un error \citep{pg-restore}.

El proceso terminó con \textbf{código de salida 0}. La figura~\ref{fig:restauracion}
muestra el comando y la finalización de la creación de las claves foráneas.

\clearpage
\subsection{Verificación de la restauración}
Se compararon los nombres de las tablas recuperadas con el índice del
respaldo. También se consultaron las secuencias, la extensión instalada
y las restricciones de integridad. El conteo de filas se realizó con
\texttt{count(*) FROM ONLY}, para evitar duplicidades por herencia.

\begin{table}[h!]
\centering
\caption{Objetos y datos verificados en la base \texttt{ejemplo}.}
\label{tab:restauracion}
\begingroup\renewcommand{\arraystretch}{1.45}
\rowcolors{2}{FondoTabla}{white}
\begin{tabularx}{.94\linewidth}{@{}Jr@{}}
\toprule\rowcolor{white}\textbf{Elemento} & \textbf{Resultado}\\\midrule
Tablas en el esquema \texttt{public} & 46\\
Secuencias en el esquema \texttt{public} & 31\\
Claves primarias & 22\\
Claves foráneas & 62\\
Restricciones \texttt{CHECK} & 55\\
Restricciones sin validar & 0\\
Filas en las tablas restauradas & 0\\
Extensión \texttt{adminpack} & Versión 2.1\\\bottomrule
\end{tabularx}\endgroup
\end{table}

\subsubsection{Correspondencia con el contenido original}
Las 46 tablas restauradas están vacías. Para verificar que este resultado
correspondía al insumo, se contaron los registros de los bloques
\texttt{COPY} del respaldo: los 46 bloques contienen cero filas y se
incluyen 31 estados de secuencias.

Las consultas sobre \texttt{linea}, \texttt{estacion} y su unión con
\texttt{componerse} devolvieron cero filas, en concordancia con el contenido
original. Por tanto, el respaldo suministrado permite recuperar la
estructura y el estado de las secuencias, pero no incorpora registros
en las tablas.

\subsubsection{Resultado}
La restauración se completó sin errores. Se recuperaron todas las tablas
identificadas en el índice y las restricciones se encuentran validadas.
Las figuras~\ref{fig:verificacion} y~\ref{fig:consultas} presentan las
comprobaciones de objetos, integridad y contenido.

\clearpage
\begin{landscape}
\paginaevidencia{evidencias/restauracion/01_inspeccion.png}
{Inspección del respaldo y comprobación de la base de destino: formato,
versiones y ausencia de tablas antes de la restauración.}{fig:inspeccion}
\paginaevidencia{evidencias/restauracion/02_restauracion.png}
{Ejecución de \texttt{pg\_restore}. Se presentan el comando, extractos del
proceso y el código de salida 0.}{fig:restauracion}
\paginaevidencia{evidencias/restauracion/03_verificacion.png}
{Verificación de las tablas, secuencias y restricciones: 46 tablas,
31 secuencias y ninguna restricción pendiente de validación.}{fig:verificacion}
\paginaevidencia{evidencias/restauracion/04_consultas.png}
{Consultas de comprobación y contraste con el archivo original:
46 bloques de datos sin filas y 31 estados de secuencias.}{fig:consultas}
\end{landscape}
